\SourceLeaf{281}%
% A544-P0281-D01
Per hanc formulam eruta est multitudo
\[
\begin{array}{r|r||r|r||r|r}
A&&A&&A&\\\hline
100&317&1000&3149&10000&31417\\
200&633&2000&6293&20000&62845\\
300&949&3000&9425&30000&94237\\
400&1257&4000&12581&40000&125629\\
500&1581&5000&15705&50000&157093\\
600&1885&6000&18853&60000&188453\\
700&2209&7000&21993&70000&219901\\
800&2521&8000&25137&80000&251305\\
900&2821&9000&28269&90000&282697\\
1000&3149&10000&31417&100000&314197
\end{array}
\]

% A544-P0281-D02
\Article{4}
Ad propositum nostrum non requiritur determinatio exacta, sed potius indagatio expressionis, quae ad multitudinem exactam quam prope velis accedere potest, dum limites in infinitum ampliantur. Sed ante omnia quum haec aliquid vagi involvant, rem exactius explicare oportet.

% A544-P0281-D03
Supponemus itaque, functiones $P,Q,R$ etc. praeter variabiles $x,y$ implicare elementum constans $k$, ita ut singulae $P,Q,R$ etc. sint functiones homogeneae trium quantitatum $x,y,k$. Hoc pacto figura per aequationes $P=0,Q=0,R=0$ etc. determinata pendebit a $k$, ita ut valoribus diversis ipsius $k$ respondeant figurae similes et respectu initii coordinatarum similiter positae, dimensionesque lineares similes valoribus ipsius $k$, areae valoribus ipsius $kk$ proportionales erunt. Denotetur iam multitudo punctorum intra figuram per $M$, area per $V$, patetque $M$ et $V$, crescente $k$, crescere debere; crescente vero $k$ in infinitum, $M$ et $V$ ad rationem aequalitatis quam proxime velis accedent, vel si elementarem claritatem postulas, proposita quantitate quantumvis parva $\lambda$, semper assignari poterit terminus talis, ut pro quolibet valore ipsius $k$ hunc terminum superante certo $\frac MV$ iacere debeat inter $1-\lambda$ et $1+\lambda$. Secundum morem suetum hoc ita indicare licet: fieri $M=V$ pro valore infinito ipsius $k$.

% A544-P0281-D04
In exemplo nostro conditio requisita locum tenet, statuendo $k=\surd A$, curvaque fit circulus, cuius area $=\pi A$, denotante $\pi$ semicircumferentiam circuli pro radio $=1$. Numeri supra traditi convergentiam luculenter addigitant.

\SourceLeaf{282}%
% A544-P0282-D01
Ceterum si operae pretium esset, facile demonstrationem illius theorematis antiquo rigore absolvere possemus, quam tamen hocce quidem loco supprimere maluimus ad difficiliora properantes.

% A544-P0282-D02
\Article{5}
In hacce commentatione limes per \emph{unicam} aequationem talem exprimetur $axx+2bxy+cyy=A$, ita quidem ut $a,b,c$ sint integri, atque $bb-ac$ numerus negativus quem statuemus $=-D$. Manifesto curva figuram definiens erit ellipsis, patetque facile, quadrata semiaxium esse radices aequationis
\[
(ac-bb)qq-(a+c)Aq+AA=0\quad\text{sive}\quad
=A\left(\frac{a+c\pm\surd(4bb+(a-c)^2)}{2(ac-bb)}\right)
\]

% A544-P0282-D03
Productum harum radicum fit $\frac{AA}{ac-bb}=\frac{AA}D$, proin area ellipsis $=\frac{\pi A}{\surd D}$. Hinc itaque colligitur, multitudinem omnium combinationum valorum integrorum ipsarum $x,y$, pro quibus $axx+2bxy+cyy$ valorem $A$ non superet, crescente $A$ continuo magis appropinquare ad $\frac{\pi A}{\surd D}$, et pro $A$ infinito huic valorem aequalem statui debere. Ceterum manifestum est, hocce respectu nihil interesse, utrum combinatio $x=0,y=0$ reliquis annumeretur, an inde excludatur. Hoc itaque modo multitudo quaesita (in ratione posteriori) nihil aliud est, nisi aggregatum multitudinum repraesentationum singulorum numerorum $1,2,3\ldots A$ per formam binariam secundi gradus $axx+2bxy+cyy$; et quum inter illos numeros alii omnino per hanc formam repraesentari nequeant, alii plures, alii pauciores repraesentationes admittant, quantitas $\frac\pi{\surd D}$ consideranda erit tamquam valor medius multitudinis repraesentationum numeri positivi indefiniti per formam quamlibet, cuius determinans $=-D$.

% A544-P0282-D04
\Article{6}
Antequam quae hinc sequantur generaliter perscrutemur, ut modus argumentationis facilius penetrari possit, casus quosdam singulares evolvere visum est. Resumamus itaque primo formam $xx+yy$, pro qua itaque multitudo repraesentationum numeri indefiniti valorem medium $=\pi$ nanciscitur. Multitudo vero repraesentationum actualium numeri dati haud difficile e principiis generalibus in Disquisitionibus Arithmeticis stabilitis determinatur. Designemus per $fA$ multitudinem repraesentationum numeri $A$, quae erit $=4$, si $A=1$ vel $2$ vel potestas binarii; $=8$, si $A$ est numerus primus formae $4n+1$, vel productum
\SourceLeaf{283}%
% A544-P0283-D01
talis numeri primi in potestatem binarii; $=0$, si $A$ est numerus primus formae $4n+3$, vel per talem numerum primum divisibilis, neque vero per ipsius quadratum; denique \emph{generaliter}
\[
\begin{aligned}
\text{vel}\quad&=4(\alpha+1)(\beta+1)(\gamma+1)\ldots\\
\text{vel}\quad&=0
\end{aligned}
\]
prout, reducto numero $A$ ad formam $2^\mu Sa^\alpha b^\beta c^\gamma\ldots$, designantibus $a,b,c$ etc. numeros primos inaequales formae $4n+1$, $S$ autem productum e numeris primis formae $4n+3$, si qui inter factores numeri $A$ semel pluriesve occurrunt, numerus $S$ est vel quadratum vel non quadratum. Patet itaque, $fA$ unice pendere a modo, quo numeri primi $3,5,7,11,13$ etc. inter factores numeri $A$ reperiuntur, ita ut generaliter statuere oporteat
\[
fA=4(3)\pdot(5)\pdot(7)\pdot(11)\pdot(13)\dotsfour
\]
si valores characterum $(3),(5),(7)$ etc. ita acceptos supponimus, ut denotante $p$ numerum primum sit

% A544-P0283-D02
\begin{description}
\item[primo] $(p)=1$, si $p$ ipsum $A$ non metitur
\item[secundo] $(p)=\alpha+1$, si $p$ est formae $4n+1$, atque $p^\alpha$ potestas summa ipsum $A$ metiens
\item[tertio] $(p)=0$, si $p$ est formae $4n+3$, atque exponens potestatis altissimae ipsius $p$ ipsum $A$ metientis est impar: denique
\item[quarto] $(p)=1$, si $p$ est formae $4n+3$, atque exponens potestatis summae ipsius $p$ ipsum $A$ metientis est par.
\end{description}
Manifesto casus primus sub secundo et quarto continetur.

% A544-P0283-D03
Hoc itaque modo termini progressionis $f1,f2,f3,f4$ etc. valde irregulariter procedunt, etiamsi quo maior multitudo sumatur, eo accuratius valor medius $=\pi$ inde surgere debeat. Aggregatum $f1+f2+f3+\ldots+fA$ denotabimus per $FA$.

% A544-P0283-D04
\Article{7}
Statuamus iam generaliter $fm+f3m=f'm$, perspicieturque facile, fieri
\[
f'A=4(5)\pdot(7)\pdot(11)\pdot(13)\ldots
\]
\SourceLeaf{284}%
% A544-P0284-D01
i. e. $f'A$ a relatione ipsius $A$ ad divisorem $3$ erit independens, unde seriei $f'1,f'2,f'3,f'4,f'5,f'6$ etc. irregularitas tum serius incipiet tum longe minor erit. Porro si statuimus
\[
f'1+f'2+f'3+f'4+\text{etc.}+f'm=F'm
\]
erit
\[
\begin{aligned}
F'3A&=F3A+f3+f6+f9+\ldots+f9A\\
&=F3A+FA
\end{aligned}
\]
Hinc facile concluditur crescente $A$ in infinitum, statui debere
\[
F'3A=4\pi A
\]
sive valorem medium terminorum seriei $f'1,f'2,f'3,f'4$ etc. esse
\[
=\tfrac43\pi
\]

% A544-P0284-D02
Simili modo statuendo generaliter $-f'm+f'5m=f''m$, fiet
\[
f''A=4(7)(11)(13)\ldots
\]
sive e serie nova $f''1,f''2$ etc. abeunt vacillationes a relatione ad numerum $5$ pendentes. Statuendoque aggregatum
\[
f''1+f''2+f''3+\ldots+f''m=F''m
\]
fiet
\[
F''5m=-F'm+F'5m
\]
unde concluditur crescente $m$ in infinitum, statui debere
\[
F''5m=\tfrac43\pi\pdot4m
\]
sive valorem medium terminorum seriei esse $=\tfrac43\pdot\tfrac45\pi$.

% A544-P0284-D03
Si eodem modo ulterius procedimus, progressiones novas formando, dum deinceps factores $(7),(11),(13),(17)$ etc. tollimus, hae continuo magis ad invariabilitatem appropinquabunt, valoresque medii deinceps novos factores $\tfrac87,\tfrac{12}{11},\tfrac{12}{13},\tfrac{16}{17},\tfrac{20}{19}$ etc. nanciscentur, ubi denominatores erunt numeri primi serie naturali, numeratores vero unitate vel maiores vel minores, prout illi sunt formae $4n-1$, vel ${}^{.}4n+1$. Quare quum hoc processu in infinitum continuato valor 
\SourceLeaf{285}%
% A544-P0285-D01
constans $4$ valori medio continuo propior fieri debeat, habemus
\[
4=\pi\pdot\tfrac43\pdot\tfrac45\pdot\tfrac87\pdot\tfrac{12}{11}\pdot\tfrac{12}{13}\ldots\text{ in inf.}
\]
sive
\[
\pi=4\pdot\frac3{3+1}\pdot\frac5{5-1}\pdot\frac7{7+1}\pdot\frac{11}{11+1}\pdot\frac{13}{13-1}\ldots
\]
Si singulae fractiones evolvuntur in series infinitas
\[
\begin{aligned}
\frac3{3+1}&=1-\frac13+\frac19-\frac1{27}+\ldots\\
\frac5{5-1}&=1+\frac15+\frac1{25}+\frac1{125}+\ldots\\
\frac7{7+1}&=1-\frac17+\frac1{49}-\frac1{343}+\ldots\\
&\text{etc.}
\end{aligned}
\]
productum facile evolvitur in
\[
1-\tfrac13+\tfrac15-\tfrac17+\tfrac19-\tfrac1{11}+\tfrac1{13}\dotsfour
\]
cuius seriei summam esse $=\tfrac14\pi$ vulgo notum est. Revera via inversa olim iam hinc aequalitas inter $\tfrac14\pi$ et productum infinitum $\tfrac43\pdot\tfrac45\pdot\tfrac87\ldotp\ldotp$ ab ill. \textsc{Euler} erutum fuerat (\emph{Introd. in analys. inf.} T. I. Cap. xv. art. 285).

% A544-P0285-D02
\Article{8}
Consideremus secundo loco formam $xx+2yy$, pro qua multitudo repraesentationum numeri indefiniti valorem medium $=\frac\pi{\surd2}$ habebit. Designando per $fA$ multitudinem repraesentationum numeri dati $A$ per istam formam, haec erit $=2$ pro $A=1$ vel $A=2$, vel quoties $A$ est potestas binarii; porro $fA=4$, quoties $A$ est aliquis e serie numerorum primorum, quorum residuum quadraticum est $-2$, sive qui sunt formae $8k+1,8k+3$, puta $A=3,11,17,19,41,43$ etc.; denique $fA=0$, quoties $A$ est numerus primus, cuius non-residuum quadraticum est $-2$, puta e serie $5,7,13,23,29,31$ etc. sive vel formae $8k+5$, vel formae $8k+7$. Generaliter vero statui debet
\[
\begin{aligned}
\text{vel}\quad fA&=2(\alpha+1)(\beta+1)(\gamma+1)\ldots\\
\text{vel}\quad fA&=0
\end{aligned}
\]
prout reducto numero $A$ ad formam $2^\mu Sa^\alpha b^\beta c^\gamma\ldots$, designantibus $a,b,c$ etc. numeros primos inaequales formae $8k+1,8k+3$, contra $S$ productum e 
